% ======================================================================
% SECTION 6: CONCLUSIONS
% Four thematic blocks: headline artifact counts, persistent themes,
% implications for patient safety and effectiveness, forward path.
% Location: 2030-gbm-1min/paper/full-paper/sections/conclusions.tex
% ======================================================================

The 2030 GBM 1-minute trial package documents one continuous
workflow across three artifact tiers: the 12 hand-authored
instruction files at
kevinkawchak/physical-ai-oncology-trials/competitions/instructions/one_minute_variant/
(approximately 130 KB total)
\cite{one-minute-variant-instructions}, the generated
kevinkawchak/robotic-surgeries/2030-gbm-1min/ tree (approximately
8.2 MB committed plus 1.5 MB of fixed overhead inside the 10 MB
GitHub cap) \cite{repo-robotic-surgeries}, and the end-to-end
execution outputs at
kevinkawchak/robotic-surgeries/2030-gbm-1min/outputs/, including the
54 column by 1001 row sensor sample table at
outputs/sensors/sensor_sample_4arm.csv that no human team could
produce within the time budget of this paper. The headline artifact
counts are: 12 instruction Markdown files, 1 2030-gbm-1min tree, 16
deterministic iterations producing 80 per-iteration Parquet files
(16 L1 50 ms plus 16 L2 1 s plus 16 L3 per-phase plus 16 events)
plus 16 L0 Zenodo pointer JSON files, 1 sensor sample table at 54
columns by 1001 rows, 240 per-arm xyz commands across 4 files, 6
curated ASCII diagrams at outputs/diagrams/, 4 ASCII bar / histogram
charts at outputs/viz/, 4 narrative reports at outputs/reports/, and
2 LLM tournaments at outputs/comparison/ and
outputs/comparison_robot_vs_human/.

Three persistent themes recur across the artifact set.
\textbf{Theme 1.} The on-premises LLM with a 1M token context window
holds every instruction file, every kinematics and schema
specification, every iteration aggregate, and every comparison
rationale in a single inference pass \cite{claude-opus-47,
claude-code}. The 1M context window is the property that makes the
end-to-end pipeline tractable from a single agent rather than
requiring brittle chained calls. \textbf{Theme 2.} The deterministic
seed 20260510 plus the frozen weighted composite score formula
(quality 0.40, time 0.25, cost 0.20, safety 0.10, patient experience
0.05) make the entire pipeline bit reproducible for fixed seed plus
parameter tuple; the run-summary report at
outputs/reports/run_summary.md preserves the exact seed and weight
record. \textbf{Theme 3.} The hourly commit cadence and the
structured artifact bundle (Parquet plus JSON plus ASCII plus PDF)
align with the FDA Real-Time Clinical Trials (RTCT) vision
\cite{fda2026realtime} that the robotic-surgeries pipeline extends
from pharmacology into the surgical theater.

The implications for patient safety and effectiveness are direct.
The on-premises LLM single-robot-error-minimization thesis is
operational through the 1 kHz heartbeat broadcast bus
\cite{iec-80601-2-77, iec-62304} and the cumulative tip force cap
of 12 N on the patient frame across the 4 cooperating arms. The
4-arm cross-check matrix means no single arm can drive the patient
force envelope outside the allowed band; an LLM-emitted xyz command
that violates the cumulative cap is rejected at the heartbeat
boundary and triggers the 5 ms E-stop with the 100 microsecond
emergency arm-park. The patient experience gain is the 240x to 480x
compression of operative duration (4 to 8 hours compressed to 60
seconds), and the safety gain is the deterministic cross-arm cap
that no single-arm system provides. The work is framed respectfully
with the FDA: it extends the RTCT proof-of-concept program from
pharmacology to the surgical theater under the Software as a Medical
Device framework \cite{fda-samd}, and it positions the United States
to remain Number 1 in the world regarding patient safety, efficacy,
and speed benefits in oncological robotic surgeries in clinical
trials.

The forward path is to bridge the synthetic to real gap. TRIPOD+AI
and CREMLS \cite{Collins2024TRIPODAI, ElEmam2024CREMLS} set the
reporting floor; \S\ref{sec:limitations} carries the formal
accounting of approximations vs generated vs executed, the 60-minute
vs 1-minute deltas, the cross-simulation limitations, and the
Track A vs Track B future directions. The Overleaf-ready bundle at
\path{2030-gbm-1min/paper/full-paper/LaTeX Source Files.zip}
uploads in one step; the GitHub commit log of this pull request
provides the audit trail for the template-population pass. The
upstream artifact chain runs through
\cite{kawchak_2026_20113157, repo-robotic-surgeries,
kawchak_2026_19994945, kawchak_2025_17774560,
kawchak_2025_15549831, kawchak_2025_17614396}. The 60-second target
is a checking step the industry no longer has to take, but the
practical insight on how close the application is to real life
today is that a real NeuroSpeed-class robot, a real patient cohort,
and a TRIPOD+AI plus CREMLS aligned reporting layer are required
before any clinical use.
